\documentstyle[%


righttag%


,11pt%


,amssymb%


,russian%               remove in Eng.


,izvru%                 Set izven% in Eng.


]{amsart}





%





\newcommand{\mes}{\operatorname{mes}}


\newcommand{\vraisup}{\operatorname{vrai\;sup}}


\newcommand{\tts}[1]{}





%\hoffset=-15mm%                  For Eng.


\setcounter{page}{32}


\god{2001}


\nomer{12~(475)} %                        Remove in Eng.


%\nomer{Vol.\,45, No.\,12}%               Set in Eng.


%\str{\pageref{begin}--\pageref{end}}%  Set in Eng.


\udk{517.977}





\author{\it В.А.\,Дыхта, Н.В.\,Деренко}


% NB:   ^^ remove in Eng.


\title{Численные методы решения задач оптимального


импульсного управления, основанные на вариационном принципе максимума}


\thanks{Работа выполнена при финансовой поддержке


Российского фонда фундаментальных исследований, грант \No\,01-01-00869.}





\date{}


%\pagestyle{myheadings}%         Set in Eng.


\pagestyle{plain} %              Remove in Eng.


\begin{document}


%\thispagestyle{plain}%          Set in Eng.


\maketitle


\label{begin}








\section{Введение}





В качественной теории 


оптимизации динамических систем с разрывными


траекториями и импульсными управлениями достигнут заметный прогресс: 


получены вариационный принцип максимума и обобщенные условия стационарности


для импульсных процессов с траекториями неограниченной вариации [1], [2],


различные варианты принципа максимума для импульсных процессов с траекториями 


ограниченной вариации [3]--[5].


В то же время наметился определенный отрыв качественной теории от 


конструктивных численных методов решения  задач нелинейного


импульсного управления. Публикации в этой области единичны [6], [7], 


носят частный характер и слабо связаны с условиями оптимальности 


импульсных процессов того или иного класса (для процессов с траекториями


ограниченной вариации авторам вообще неизвестны какие-либо работы


по численным методам).





В данной статье предлагаются методы решения нелинейных задач


импульсного управления с траекториями неограниченной вариации


(класса $L_\infty$), которые базируются на вариационном принципе 


максимума и обобщенном условии стационарности [1], [5]. 


Техника, развитая в [1], позволяет предложить численные


методы, конструкции которых имеют естественный и компактный


вид. Для построения методов градиентного типа она была


применена в [8] в случае менее общей задачи, нежели 


рассматриваемая ниже.





\section{Постановка задачи}





Пусть $T=[t_0,t_1]$ --- фиксированный отрезок времени,


$\cal W$ --- класс $m$-мерных измеримых ограниченных 


вектор-функций, постоянных при $t<t_0$ и $t>t_1$,


$\cal D$ --- оператор обобщенного дифференцирования.





Будем рассматривать задачу оптимального 


управления, содержащую так называемые ограничения на текущие


значения импульса,


\begin{gather}


J(w)=l(x(t_1+),w(t_1+)) \to \inf,\\


{\cal D} x=f(t,x,w)+G(t,x,w){\cal D}w,\\


x(t_0-)=x_0, \quad w(t_0-)=w_0,\\


w(t) \in W.


\end{gather}


Здесь $l(x,w)$ --- скалярная функция, $x(t)\in R^n$, $x(t\pm)$ ---


односторонние пределы функции $x$ в точке $t$, уравнение


(2) трактуется в смысле распределений [2], $W \subset R^m$ --- 


заданное множество, управление $w \in {\cal W}$.





Предполагаются выполненными следующие условия:





1) функция $l(x,w)$ непрерывно дифференцируема;





2) вектор-функция $f(t,x,w)$ и ее производные по $x$ и $w$


непрерывны;





3) матричная функция $G(t,x,w)$ дважды непрерывно дифференцируема;





4) выполнено условие Фробениуса для матрицы $G$ по переменным $(x,w)$


$$


G_{ix}(t,x,w) G_j(t,x,w)+G_{iw_j}(t,x,w)=


G_{jx}(t,x,w) G_i(t,x,w)+G_{jw_i}(t,x,w), \quad i,j=1,2,\dotsc,m,


$$


где $G_1,\dotsc, G_m$ --- столбцы матрицы $G$, а также удовлетворяется 


условие роста $| G(t,x,w)| \leq c(1+|x| +|w|)$ при некотором


$c>0$.


 


Заметим, что предположение 4) необходимо для возможности 


корректного определения решения уравнения в распределениях 


(2) [2]. Это уравнение является расширением обычной


системы уравнений


\begin{equation}


\Dot{x}=f(t,x,w)+G(t,x,w)u,\quad \Dot{w}=u


\end{equation}


с управлением $u(\cdot)$ класса $L_\infty$ и липшицевыми 


траекториями. При этом функции $x(\cdot)$, $w(\cdot)$, 


связанные уравнением в распределениях (2), 


являются пределом почти всюду сходящейся последовательности 


траекторий системы (5), равномерно ограниченной в


$C$ и обладающей свойством $u_n \to {\cal D}w$


в смысле обобщенных функций [4], [5]. Поэтому при выпуклом 


множестве $W$ переход от системы (5) к уравнению (2) не


ведет к уменьшению инфимума функционала $J$ (это верно лишь


для рассматриваемой базовой задачи со свободным правым 


концом). В задаче (1)--(4) управлением будем считать функции


$w\in {\cal W}$ со значениями в $W$, которые порождают импульсные


воздействия вида $u={\cal D} w$.





\section{Сведение задачи импульсного управления к обычной}





В основе предлагаемых методов решения задачи импульсного 


управления (1)--(4) лежит возможность преобразования ее к


классической задаче оптимального управления, к которой


применимы известные методы, с последующим представлением


этих методов непосредственно в терминах рассматриваемой


задачи.





Сведение задачи импульсного управления к обычной 


задается заменой фазовых координат по формуле [1], [2], [5]


\begin{equation}


y(t)=\xi(t,x(t),w_0,w(t)),


\end{equation}


где функция $\xi(t,y,w,w_0)$ --- решение вполне интегрируемой


(в силу предположения 4)) системы в частных производных


\begin{equation}


\frac{\partial \xi}{\partial w} = G(t, \xi,w), \quad  \xi|_{w=w_0}=y,


\end{equation}


а $x(\cdot)$ --- траектория уравнения (2), соответствующая 


допустимому управлению $w(\cdot)\in{\cal W}$. Если для каждого такого


управления положить $h=w(t_1+)$, то преобразованием (6)


задача (1)--(4) редуцируется к следующей:


\begin{gather}


I(w,h)={\cal L}(y(t_1),h) \to \inf,\\


\Dot{y}=g(t,y,w), \quad y(t_0)=x_0,\\


w(t)\in W, \quad h\in W.


\end{gather}


Здесь ${\cal L}(y,h)=l(\xi(t_1,y,h,w_0),h)$, 


$g(t,y,w)=\xi_t(t,x,w_0,w)+\xi_x(t,x,w_0,w)


f(t,x,w)|_{x=\xi(t,y,w,w_0)}$;


$w$~--- измеримое ограниченное 


управление, $h$ --- оптимизируемый параметр, $y$ --- липшицевая 


траектория, связанная с соответствующей траекторией уравнения


(2) формулой обращения для (6)


\begin{equation}


\begin{array}{c}


x(t)=\xi(t,y(t),w(t),w_0), \ \ t\in(t_0,t_1),\\[2mm]


x(t_0-)=x_0,\ \ x(t_1+)=\xi(t_1,y(t_1),h,w_0).


\end{array}


\end{equation}


 


Переход от исходной задачи (1)--(4) к редуцированной


(8)--(10) в принципе можно использовать для анализа 


всевозможных свойств задачи импульсного управления и, в 


частности, для ее непосредственного решения. Однако мы будем


исходить из предположения, что практически решить в 


аналитическом виде систему (7) невозможно, так что описанная


редукция лишь теоретически возможна, а конструктивные 


результаты для исходной задачи должны получаться расшифровкой


соответствующих аналогов для задачи (8)--(10).





\section{Основная формула приращения функционала\protect\\ и вариационный 


принцип максимума}





Введем сопряженные переменные


$\psi$, $\psi_w$, отвечающие переменным $x$, $w$ 


соответственно, а также сопряженную систему в распределениях для


исходной задачи


\begin{equation}


\begin{array}{c}


{\cal D} \psi = -H_x(t,x,w,\psi,{\cal D} w), \quad


{\cal D} \psi_w = -H_w(t,x,w,\psi,{\cal D} w), \\[2mm]


\psi(t_1+)=-l_x(x(t_1+),w(t_1+)), \quad \psi_w(t_1+)=-l_w(x(t_1+),w(t_1+)).


\end{array}


\end{equation}


Здесь $H(t,x,w,\psi,u)= \langle


\psi,f(t,x,w)+G(t,x,w)u\rangle$ --- функция Понтрягина исходной задачи.





Пусть $w(\cdot)$ --- некоторое допустимое управление, $x(\cdot)$ ---


соответствующая траектория уравнения (2), $\psi(\cdot)$, $\psi_w(\cdot)$ ---


решение системы (12) и $\omega(\cdot)$ --- другое допустимое 


управление в исходной задаче. Для $\tau \in[0,1]$ положим


$$


{\mathrm w}_\omega(t,\tau)=w(t)+\tau(\omega(t)-w(t)), \quad t\in T,


$$


и введем так называемые предельные дифференциальные системы


для уравнений (2), (12) вдоль импульсного процесса $(x(\cdot),w(\cdot))$


\begin{equation}


\begin{array}{c}


z'(\tau)=G(t,z(\tau),{\mathrm w}_\omega(t,\tau))(\omega(t)-w(t)),\quad


q'(\tau)=-H_{xu}(t,z(\tau),{\mathrm w}_\omega


(t,\tau),q(\tau))(\omega(t)-w(t)), \\[2mm]


z(0)=x(t),\quad q(0)=\psi(t),


\end{array}


\end{equation}


где штрих означает дифференцирование по $\tau$.





В формулу приращения функционала $J$ войдет некоторая


функция от $\omega$ и решений $z(\tau)$, $q(\tau)$ системы (13), 


а точнее --- зависящий от $t\in T$, как параметра, функционал.


Чтобы определить его, положим


\begin{equation}


P(t,x,w,\psi)=H_w(t,x,w,\psi)-\frac{d}{dt}H_u(t,x,w,\psi),


\end{equation}


где полная производная по $t$ вычисляется по обычному 


правилу в силу уравнений (2), (12) и связи $\Dot{w}=u$ 


(от вектора $u$ эта производная не будет зависеть из-за 


предположения 4)). Пусть


\begin{equation}


{\cal P} (\omega(t),w(t),t)=\Big\langle \int_0^1P(t,z(\tau),{\mathrm w}_\omega


(t,\tau),q(\tau))d\tau, \omega(t)-w(t)\Big\rangle,


\end{equation}


где $z(\tau)$, $q(\tau)$ --- зависящие от $\omega(t)$ решения системы (13).





Для построения методов важной является





\begin{thm}


Приращение функционала $J(w)$ может быть


представлено в виде


\begin{multline}


\Delta_\omega J(w):=J(\omega)-J(w)=-\int_T{\cal P}(\omega(t),w(t),t)dt-\\


-\langle H_u(t_1,x(t_1+),w(t_1+),


\psi(t_1+))+\psi_w(t_1+),\omega(t_1+)-w(t_1+)\rangle + \\


+ o(\| \omega - w\|_1 +|\omega(t_1+)-w(t_1+)|),


\end{multline}


где интегрант ${\cal P}$ определен равенством $(15)$, а символ $\|\cdot\|_1$


означает норму в пространстве~$L_1$.


\end{thm}





\begin{pf}


В редуцированной задаче


(8)--(10) справедливо известное разложение функционала


\begin{multline}


\Delta_{\omega, \gamma} I(w,h):=I(\omega,\gamma)-I(w,h) = \\


=-\int_T \Delta_\omega{\cal H}(t,y(t),p(t),w(t))dt+


\langle {\cal L}_h(y(t_1),h),\gamma-h\rangle + 


o(\|\omega-w\|_1 + |\gamma - h|).


\end{multline}


Здесь ${\cal H}(t,y,p,w)=\langle p,g(t,y,w)\rangle $ ---


функция Понтрягина задачи (8)--(10), $\Delta_\omega {\cal H}$


--- ее частное приращение по управлению,


$y(\cdot)$ --- траектория задачи (8)--(10), 


соответствующая фиксированному управлению


$w\in L_\infty(T,W)$, $p(\cdot)$ --- решение сопряженной системы 


редуцированной задачи, т.\,е.


$$


\Dot{p} =-{\cal H}_y(t,y,p,w), \quad p(t_1)=-{\cal L}_y(y(t_1),h),


$$


$\gamma=\omega(t_1+)$, $h=w(t_1+)$ --- сравниваемые значения параметров.





Используя подходящую модификацию рассуждений из ([1], леммы 3, 5), 


нетрудно убедиться, что


$$


\Delta_\omega {\cal H}(t,y(t),p(t),w(t))={\cal P}(\omega(t),w(t),t),


$$


так что интегранты в (16) и (17) совпадают. Равенство 


вторых слагаемых в этих разложениях следует из определения


функции $\cal L$, равенств (11) и условий трансверсальности (см. (12)). 


Поскольку функционалы $J$ и $I$ на соответствующих процессах 


совпадают, то тождественность разложений (16), (17)


установлена.


\end{pf}





Из разложения (17) с учетом произвольности выбора


управления $\omega\in L_\infty(T,W)$ вытекает следующее необходимое


условие оптимальности для задачи импульсного управления --


вариационный принцип максимума (ВПМ) в интегральной форме.


\medskip





{\bf Следствие (ВПМ).} Если допустимое управление $w(t)$


оптимально, то выполняется неравенство


\begin{equation}


\delta_\omega J(w):=-\int_T{\cal P}(\omega(t),w(t),t)dt \ge 0 


\ \; \forall\, \omega\in L_\infty(T,W).


\end{equation}


Кроме того, если предел $w(t_1-)$ существует, то оптимальное


значение $\ov{h} =w(t_1+)$ является решением экстремальной задачи


\begin{equation}


\begin{array}{c}


l(z(1),h) \to \min, \quad


z'=G(t_1,z,w(t_1-)+\tau(h-w(t_1-)))(h-w(t_1-)), \\[2mm]


z(0)=x(t_1-), \quad h\in W.


\end{array}


\end{equation}





Заметим, что экстремальное условие (19), 


характеризующее оптимальность терминального импульса, следует из того,


что параметр $h$ входит лишь в функционал задачи (8)--(10),


а множество всех фазовых точек, в которые можно попасть из


положения $x(t_1-)$ за счет импульса управления в момент $t=t_1$,


совпадает с множеством значений $z(1)$ траекторий


системы (19). При выпуклом множестве $W$ условие (19) 


может быть локализовано, т.е. ослаблено за счет рассмотрения


значений $h$, близких к $w(t_1-)$. Понятно, что условие (18)


является эквивалентом принципа максимума Понтрягина для 


редуцированной задачи (8)--(10).





\section{Схема улучшения импульсного управления}





Из условия оптимальности (18) следует, что проверка ВПМ для управления


$w(\cdot)$ сводится к решению задачи


\begin{equation}


\delta_\omega J(w) \to \min, \quad \omega \in L_\infty(T,W).


\end{equation}


Поскольку в ней допустимые функции $\omega(\cdot)$ разрывны, то эта


интегральная задача равносильна поточечной максимизации


функции $\omega \to {\cal P}(\omega, w(t),t)$


при $t\in T$ и $\omega \in W$. Формально ситуация выглядит 


аналогично имеющейся в методах, основанных на принципе максимума


Понтрягина [9]--[11] в обычных задачах управления, и любая из


версий этих методов порождает соответствующий аналог для 


решения задачи импульсного управления. Однако отличительной


особенностью рассматриваемой ситуации является, как следует


из (15), неявная зависимость функции $\cal P$ от $\omega$ 


(аналогично обстоит дело и во вспомогательной терминальной задаче


(19)). Из (13), (15) видно, что при фиксированном $t\in T$


максимизация функции $\cal P$ по $\omega \in W$ представляет собой


задачу оптимизации параметра в гамильтоновой системе (13).


Следовательно, нужно располагать достаточно эффективным


методом решения вспомогательной задачи данного класса. 


Методы градиентного типа, обсуждаемые в следующем пункте, дают


один из подходов к решению вспомогательной задачи. 


Необходимо заметить, что ситуация значительно упрощается в тех 


приложениях, в которых предельные системы (13) могут быть


проинтегрированы аналитически.





Допустим теперь, что решение $\ov \omega(\cdot)$


задачи (20) известно, и рассмотрим возможную схему улучшения управления


$w(\cdot)$, использующую, например, оптимальное игольчатое 


варьирование [9]. Заметим, что параллельно должна решаться и 


терминальная задача (19) улучшения параметра $h$.





Фиксируем параметр $\alpha \in [0,1]$, образуем множество 


варьирования $T_\alpha$, $\mes T_\alpha = \alpha(t_1-t_0)$


и рассмотрим проварьированное управление


\begin{equation}


w_\alpha (t)=w(t) + \chi_\alpha (t) (\ov \omega(t)-w(t)),


\end{equation}


где $\chi_\alpha (t)=\chi_{T_\alpha}(t)$ -- характеристическая 


функция множества $T_\alpha$ (т.\,е. $\chi_\alpha (t)=1$ при $t\in T_\alpha$


и равна нулю в противном случае). Множество $T_\alpha$ выберем 


следующим образом. Пусть


$m(t)={\cal P} (\ov \omega (t), w(t),t) \ge 0$


--- невязка выполнения ВПМ. Если $\vraisup m(t) > 0$ (т.\,е.


условие (18) не выполнено), то положим


$T_\alpha = \{ t\in T: m(t)\ge \lambda\}$,


где число $\lambda = \lambda(\alpha)$ выбирается из условия


$\int\limits_T \chi_\alpha(t)dt = \alpha(t_1-t_0)$.


Далее параметр варьирования $\alpha$ находится из условия 


уменьшения целевого функционала $J(w_\alpha)<J(w)$. При нарушении


критерия оптимальности (18) такой выбор $\alpha$ всегда


возможен, что следует из разложения (16).





Конечно, можно использовать и другие известные 


процедуры игольчатого варьирования управления, более 


простые в реализации. Все они порождают, вообще говоря, 


разрывные управления типа (21), а значит, импульсы в правой


части уравнения (2) и разрывы соответствующей траектории.





Для иллюстрации работы с ВПМ приведем





\begin{exam}


В обычной постановке он содержит медленные


фазовые переменные ($x$) и быстрые ($y$):


\begin{gather*}


\Dot x =\varphi(x,y), \ \ \Dot y=f(x,y)+u, \ \ x(0)=x_0,\ \ y(0)=y_0,\\


|y| \le 1, \ \ J=\int_0^1[d(x,y)+\langle c(y),u\rangle]dt \to \inf.


\end{gather*}


Здесь $x\in R^{n_1}$, $y\in R^{n_2}$, $u(\cdot)\in L_\infty$ --- управление,


$\varphi$, $f$, $d$, $c$ --- известные функции соответствующих размерностей.





Перейдем в этой модели к импульсной постановке, положив


$w=y$, $u= {\cal D}w-f(x,w)$. Тогда получим задачу в форме


(1)--(4), причем меньшей фазовой размерности и с 


ограничением на текущие значения импульса вместо исходных фазовых,


\begin{gather*}


\Dot x = \varphi(x,w), \ \ x(0) = x_0,\\


{\cal D} \eta = d(x,w)- \langle c(w),f(x,w) \rangle +


\langle c(w),{\cal D}w \rangle :=F(x,w,{\cal D}w),\\


\eta (0-)=0,  \ \ w(0-) = y_0,\\


|w|\le 1, \ \ J(w)=\eta (1+) \to \inf


\end{gather*}


(т.\,к. $x$ --- медленная переменная без разрывов, то для нее


записано обыкновенное дифференциальное уравнение).





Пусть $(x(\cdot),\eta(\cdot),w(\cdot))$ --- произвольный допустимый 


импульсный процесс. Тогда $\psi_\eta \equiv -1$, $\psi_x$ --- медленная 


сопряженная компонента, соответствующая $x$, и $z_x(\tau)=x(t)$,


$q_x(\tau)=\psi_x(t)$, $q_\eta(\tau) \equiv -1$, поскольку 


предельные уравнения для них тривиальны.


Далее,


\begin{gather*}


P(t,x,\eta,w)=\varphi'_w(x,w)\psi_x + [d'_w(x,w)-f'_w(x,w) c(w)-


c'_w(w) f(x,w)] \psi_\eta,\\


\frac{dz_\eta}{d\tau} =


c(w+\tau (\omega-w))(\omega-w), \ \ z_\eta(0)=\eta(t).


\end{gather*}


С учетом этих соотношений нетрудно получить


\begin{multline*}


{\cal P}(\omega,w) = \int\limits_0^1 \frac{d}{d\tau}


[\langle\psi_x(t),\varphi(x(t),{\mathrm w}_\omega(t,\tau))\rangle+\\


%


+\langle c({\mathrm w}_\omega(t,\tau)),f(x(t),{\mathrm w}_\omega (t,\tau))-


d({\mathrm w}_\omega(t,\tau))\rangle] d\tau =


\langle \psi_x,\varphi\rangle + \langle c,f \rangle - d|_{\tau=0}^{\tau=1}.


\end{multline*}


Поэтому экстремальное условие метода ВПМ (18) 


равносильно максимизации по $\omega \in [-1,1]$ функции


$$


Q(\omega,t):=\langle \psi_x (t),\varphi(x(t),\omega)\rangle +


\langle c(\omega),f(x(t),\omega)\rangle - d(x(t),\omega).


$$





Терминальная задача (19), как легко убедиться, 


сводится к следующей:


$$


C(h):=\int_{y_0}^h c(w)dw \to \min, \ \ |h| \le 1.


$$


Hи в одном из известных методов решения задач


оптимального управления подобные нелокальные экстремальные


задачи по фазовым переменным не встречаются.





Итак, в данной модели импульсного управления возможно


аналитическое определение функции $\cal P$, что снимает 


трудности реализации метода. Заметим, что в рамки этого примера


укладывается задача управления двухзвенным манипулятором в


постановке [2].


\end{exam}





\section{Алгоритмические аспекты реализации}





В реализуемых версиях методов улучшения импульсного управления естественно


ограничиться классом распределений, для элементов которого


конструктивно могут вычисляться траектории исходного и 


сопряженного уравнений в распределениях. Для этой цели введем


подмножество ${\cal W}_K \subset {\cal W}$, состоящее из 


функций $w(\cdot) \in {\cal W}$, ограничение которых на 


отрезок $T$ представляет собой кусочно-непрерывную 


функцию, липшицевую на интервалах непрерывности.


Тогда соответствующее распределение ${\cal D} w$ (импульсное 


управление) представляет собой векторную меру вида


\begin{equation}


{\cal D}w = u(t) + \sum_{i=1}^r[w(s_i)]\delta(t-s_i),


\end{equation}


где $u(t)=\Dot w(t)$, $s_i \in T$ --- точки разрыва функции $w(\cdot)$,


$[w(s_i)]=w(s_i+)-w(s_i-)$ --- ее скачки, $\delta(t-s_i)$ --- функция 


Дирака, сосредоточенная в точке $s_i$.





В случае распределений вида (22) интегрирование систем


(2), (12) особенно просто [2], [5], [8]: на интервалах 


непрерывности $w(\cdot)$ функции $x$, $\psi$, $\psi_w$ 


суть решения соответствующих обычных дифференциальных уравнений 


с управлением ${\cal D} w = u(t)$, а в точках разрыва


$w(\cdot)$ скачки решений этих уравнений вычисляются с помощью 


предельных систем. А именно, если $s_i$ ---


момент приложения импульса, то


$x(s_i+)= \wt z(1)$, $\psi(s_i+)=\wt q(1)$,


где $\wt z$, $\wt q$ --- решения системы (13) при замене


$t$, $\omega(t)$, $w(t)$, $x(t)$, $\psi(t)$ на 


$s_i$, $w(s_i+)$, $w(s_i-)$, $x(s_i-)$, $\psi(s_i-)$ 


соответственно (надо иметь в виду, что интегрирование 


сопряженных уравнений ведется в обратном времени).





Остановимся теперь на методах градиентного типа для


решения вспомогательной задачи (20). Прежде всего заметим,


что из разложения (16) следует, что градиент функционала


$\omega(\cdot) \to \delta_\omega J(w)$ совпадает с градиентом функционала 


$w \to J(w)$. 


Рассматривая разложение (16) при малых $\| \omega -w\|_\infty$, или,


что то же самое, первую вариацию функционала $\omega \to {\cal P}$ вида


(15) при малой разности $|\omega(t)-w(t)|$, получим


формулу для градиента функционала $J$ в задаче импульсного управления


\begin{equation}


\nabla J(w) = -P(t,x(t),w(t),\psi(t)).


\end{equation}


Здесь функция $P$ определена равенством (14). Bидим,


что, во-первых, вычисление этого градиента связано только с


интегрированием исходной и сопряженной систем и, во-вторых,


формула (23) существенно отличается от градиента 


функционала в обычной задаче оптимального управления, которая 


соответствует замене обобщенной дифференциальной связи (2)


обычной вида (5). Учитывая, что функционал $J$ зависит не


только от значений $w$ при $t\in T$, но и от параметра $h=w(t_1+)$,


нетрудно получить из (16) и градиент $\nabla_hJ$ по этому 


параметру, равный, очевидно, коэффициенту при разности


$(\omega(t_1+)-w(t_1+))$.





Любопытно, что обе составляющие градиентов $\nabla J(w)$ и


$\nabla_hJ(w)$ могут быть получены формальным образом из 


стандартной формулы первой вариации функционала для обычной


задачи путем интегрирования по частям.


 


Итак, в случае выпуклого множества $W$ задача 


импульсного управления (или вспомогательная задача ВПМ (20)) может


решаться применением методов проекции градиента или 


условного градиента (при $W=R^m$ --- скорейшего спуска) с 


использованием описанных конструкций $\nabla J$ и $\nabla_hJ$, задающих 


направление спуска. Более трудоемкую процедуру игольчатого


варьирования (21) целесообразно применять после серии


итераций каким-либо из градиентных методов.


\medskip





{\bf Пример 2} с разрывной оптимальной траекторией был предложен


А.А.\,Милютиным в качестве тестовой задачи [10]:


$$


J = \int_0^2x^2(t)dt \to \inf, \ \


\Dot x = u, \ \ x(0)=0, \ \ x(2)=0,


\ \ x(t) \ge 1 \ \text{ при } \ t\in [1/2,3/2].


$$


При переходе от этой задачи к импульсной постановке ($w=x$,


${\cal D} w=u$) получаем элементарную задачу


\begin{gather*}


J = \int_0^2 w^2(t)dt \to \inf, \ \ w(0-)=0, \ \ w(2+)=0,\\


w(t) \ge 1, \ \ t\in [1/2,3/2].


\end{gather*}


Bидим, что, во-первых, как и в примере 1, фазоограничения


исходной задачи перешли в ограничение на текущие


значения импульса и, во-вторых, выпуклое множество $W$ в


данном примере зависит от $t$ (что принципиального значения


не имеет).





Функция $\cal P$ при любом допустимом $w$ без труда 


находится аналитически: ${\cal P}(\omega,w)=-\omega^2 + w^2$. 


Поэтому $\ov \omega (t) = \chi_{[1/2,3/2]}(t)$


и, очевидно, при варьировании методом (21) следует принять


$\alpha =1$, $T_\alpha = T_1 = T$ для любого начального 


приближения $w(t)$. Соответствующее проварьированное управление


$w_1(t)=\ov \omega(t)$ сразу порождает оптимальное распределение


${\cal D}w_1=\delta(t-1/2)-\delta(t-3/2)$.





При численном решении примера путем комбинации методов


типа условного градиента при $t\in [1/2,3/2]$ и скорейшего спуска


при $t \notin [1/2,3/2]$ нетрудно получить управление $\wt w$,


отличающееся от оптимального в равномерной метрике не более,


чем на $10^{-3}$.





Ясно, что без учета импульсного характера оптимального


решения в исходной задаче с фазоограничениями обычными


методами подобные результаты вряд ли можно получить [10].


\medskip





{\bf Пример 3} получен путем импульсного расширения задачи,


примененной в [11] для тестирования одного из алгоритмов


второго порядка,


\begin{gather*}


J = x(1+) + w^2(1+) \to \inf,\\


{\cal D} x = (w+1)^2 + (t-1)w{\cal D}w, \ \ x(0-)=0, \ \ w(0-)=0.


\end{gather*}


Здесь, как и в примере 2, аналитически находится функция


$$


{\cal P} (\omega, w) = \bigg(-2-w-\frac{\omega -w}{2}\bigg)(\omega - w),


$$


и оптимальное управление $w_1(t) \equiv -2$ находится на первой


итерации метода. Ему соответствуют распределение ${\cal D} w_1 = -2\delta(t)


+2\delta(t-1)$ и траектория $x_1(t)=t-2$, $t\in(0,1)$, $x_1(1+)=-1$.





Предложенный метод несравненно проще в реализации, чем


метод из [11], в котором требуется аналитически решать


систему дифференциальных уравнений в частных производных


типа (7).


 


Применим теперь для решения данного примера метод градиентного типа.


 


Формулы для вычисления градиента функционала $J$ имеют


вид ($u={\cal D}w$)


\begin{gather*}


\nabla J(w)=\frac{d}{dt}H_u(t,x,w,\psi) - H_w(t,x,w,\psi)=w+2,\\


\nabla_hJ(w) = -H_u(1,x(1+),w(1+),\psi(1+)) + l_w(x(1+),w(1+)) =


2w(1+) = 2h.


\end{gather*}


Нулевому приближению $w^0(t)=1$, $t\in (0,1)$, $h^0=1$ 


соответствует разрывная траектория $x^0(t)=4t-1/2$ и $\nabla J(w^0)=3$,


$\nabla_h J(w^0) =2$. Проварьированному по методу скорейшего


спуска управлению


\begin{gather*}


w^0_\alpha(t)=w^0(t) - \alpha \nabla J(w^0) = 1-3\alpha, \ \ 


h^0_\beta = h^0-\beta \nabla_h J(w^0) = 1-2\beta, \ \ \alpha, \beta >0,


\end{gather*}


отвечает траектория


$x^0_\alpha(t)=(2-3\alpha)^2t-0,5(3\alpha-1)^2$, $t\in(0,1)$.


Минимизация по $\beta$ функционала


$J(w^0_\alpha)=x^0_\alpha(1) + {h^0_\beta}^2$ 


(т.\,е. оптимизация терминального импульса) приводит к результату


$\ov \beta = 1/2$, $\wt h = 0$. Далее, несложно найти, что


$\min\limits_{\alpha \ge 0}J(w^0_\alpha) = -1$ при $\ov \alpha =3/4$.


 


Таким образом, первая же итерация метода скорейшего


спуска приводит к импульсному режиму


\begin{gather*}


\wt w(t)=-2, \ \ \wt w(0-)=0,\ \ \wt w(1+)=\wt h = 0,\\


\wt x(t) = t-2, \ \ t\in (0,1], \ \ \wt x(0-) = 0,


\end{gather*}


который является стационарным, т.\,к. выполняются условия


$\nabla J(\wt w) \equiv 0$, $\nabla_h J(\wt w)=0$.





Рассмотренные примеры наглядно демонстрируют 


преимущества перехода к задачам в импульсной постановке и 


применения специальных методов улучшения управления, использующих


теорию импульсного управления.








\begin{thebibliography}{99}





\bibitem{M1}


Дыхта В.А. {\em Вариационный принцип максимума и квадратичные 


условия оптимальности импульсных и особых процессов} // 


Сиб. матем. журн. -- 1994. -- T.\,35. -- \No\,1. -- C.\,70--82.


\bibitem{M2}


Завалищин С.Т., Сесекин А.Н. {\em Импульсные процессы:


модели и приложения}. -- М.: Наука, 1991. -- 256 c.


\bibitem{M3}


Миллер Б.М. {\em Метод разрывной замены времени в задачах оптимального


управления импульсными и дискретно-непрерывными системами} // Автоматика 


и телемеханика. -- 1993. -- \No\,12. -- С.\,3--32.


\bibitem{M4}


Орлов Ю.В. {\em Теория оптимальных систем с обобщенными


управлениями}. -- М.: Наука, 1988. -- 187 с.


\bibitem{M5}


Дыхта В.А., Самсонюк О.Н. {\em Оптимальное импульсное управление 


с приложениями}. -- М.: Физматлит, 2000. -- 256 с.


\bibitem{M6}


Орлов Ю.В., Разумовский Д.Д. {\em Численные методы 


решения задач оптимального обобщенного управления} // Автоматика и


телемеханика. -- 1993. -- \No\,5. -- C.\,44--51.


\bibitem{M7}


Батурин В.А., Урбанович Д.Е. {\em Приближенные методы 


оптимального управления, основанные на принципе расширения}. --


Новосибирск: Наука. Сиб. предприятие РАН. -- 1997. -- 175~с.


\bibitem{M8}


Дыхта В.А., Деренко Н.В. {\em Численные методы решения


задач импульсного управления, основанные на обобщенном условии 


стационарности} // Сб. тр. Всероссийск. школы


``Компьютерная логика, алгебра и интеллектное управление.


Проблемы анализа устойчивости развития и стратегической


стабильности''. -- Иркутск: Изд-во ИрВЦ СО РАН, 1994. -- Т.\,2. -- C.\,59--70.


\bibitem{M9}


Срочко В.А.


{\em Итерационные методы решения задач оптимального 


управления}. -- М.: Физматлит, 2000. -- 160 с.


\bibitem{M10}


Федоренко Р.П. {\em Приближенное решение задач 


оптимального управления}. -- М.: Наука, 1978. -- 488 c.


\bibitem{M11}


Гурман В.И., Батурин В.А., Москаленко А.И. и др.


{\em Методы улучшения в вычислительном эксперименте}. -- Новосибирск:


Наука. Сиб. отд., 1988. -- 184~c.





\end{thebibliography}





\vskip 2cm





\post{\it Иркутская государственная}{\it Поступила}


\post{\it экономическая академия}{03.09.2001}





\label{end}


\end{document}





